\slajdid{08-s003}
\begin{frame}[label=08-s003]{REALNA LOGIČKA KOLA}
\gusttekst\setlength{\parskip}{1pt}
\begin{center}\begin{tikzpicture}[circuit logic US,DEoznaka,baseline=(g.center)]
\node[not gate US,draw](g){};\draw(g.input)--++(-.5,0)node[left]{$a$};\draw(g.output)--++(.5,0)node[right]{$b$};\end{tikzpicture}\quad $a$ i $b$ signali\end{center}
$a$ je signal koji logičko kolo tumači, shvata, kao logičku jedinicu ili nuli\par
$b$ je signal koji logičko kolo generiše, postavlja, daje kao logičku jedinicu ili nuli\par
Logička funkcija kola je\quad $B=\overline A$\quad gde je $B$ logičko stanje signala $b$ i $A$ logičko stanje signala $a$\par
Funkcija kola je\quad $b=f(a)$\quad i po pravilu je nelinerana\par
\begin{columns}[c,onlytextwidth]
\column{.54\textwidth}\centering\begin{tikzpicture}[x=1cm,y=.75cm,font=\sitneoznake,>=Latex]
\fill[DEcrvena!12](1,1.25)rectangle(2.3,1.8);\fill[DEcrvena!12](3.2,0)rectangle(3.7,.55);
\foreach \y/\lab in {1.25/a,0/b}{\draw[->](0,\y-.1)--(0,\y+.95)node[left]{$\lab$};\draw[->](-.1,\y)--(5.8,\y)node[below]{$t$};}
\draw[thick](0,1.8)--(1,1.8)--(2.3,1.25)--(5.5,1.25);
\draw[thick](0,0)--(3.2,0)--(3.7,.55)--(5.5,.55);
\node at(.45,1.5){1};\node at(4.2,1.5){0};\node at(.45,.3){0};\node at(4.2,.3){1};
\node[anchor=west](sh)at(2.6,2.15){Šta shvata logičko kolo};\draw[->](sh.west)--(2,1.7);
\draw[->,DEcrvena](1.55,1.56)to[bend right=45](3.3,.3);
\node[anchor=north](da)at(3.7,-.26){Šta i kada daje logičko kolo};\draw[->](da.north)--(3.48,.18);
\end{tikzpicture}

\column{.43\textwidth}\centering
\begin{tikzpicture}[circuit logic US,DEoznaka]
\node[not gate US,draw](g){};\draw(g.input)--++(-.6,0)node[left]{$a_i(t)$};\draw(g.output)--++(.6,0)node[right]{$b_o(t)$};\end{tikzpicture}\par
Indeks $i$ – input, ulaz, fizička veličina na ulazu koja nosi informaciju o logičkim nulama i jedinicama\par
Indeks $o$ – output, izlaz, fizička veličina na izlazu koja nosi informaciju o logičkim nulama i jedinicama
\end{columns}\medskip
Signal ne može trenutno da promeni svoje stanje.\par
Mora da prođe sve vrednosti između dva stanja, makar smatrali da je taj vremenski interval beskonačno kratak.\par
\centering\alert{Signal je fizička veličina koja nosi informacije.}
\end{frame}
